\documentclass[11pt,a4paper]{article}

% ---- Encoding & fonts ----
\usepackage[utf8]{inputenc}
\usepackage[T1]{fontenc}
\usepackage{lmodern}
\usepackage{microtype}

% ---- Layout ----
\usepackage[a4paper,margin=1in]{geometry}
\usepackage{setspace}
\usepackage{parskip}
\usepackage{enumitem}

% ---- Math & symbols ----
\usepackage{amsmath,amssymb,amsthm}
\usepackage{mathtools}

% ---- Tables ----
\usepackage{booktabs}
\usepackage{array}
\usepackage{tabularx}

% ---- Code & algorithms ----
\usepackage{listings}
\usepackage{algorithm}
\usepackage{algpseudocode}
\usepackage{verbatim}

% ---- Graphics & color ----
\usepackage{xcolor}
\usepackage{graphicx}

% ---- Hyperlinks (load last among common packages) ----
\usepackage[hidelinks,colorlinks=false,bookmarks=true]{hyperref}

% ---- listings style ----
\definecolor{codebg}{rgb}{0.97,0.97,0.97}
\definecolor{codeframe}{rgb}{0.80,0.80,0.80}
\definecolor{codekey}{rgb}{0.10,0.30,0.60}
\definecolor{codecom}{rgb}{0.30,0.55,0.30}

\lstdefinelanguage{json}{
  morestring=[b]",
  morestring=[d]',
  morecomment=[l]{//},
}

\lstset{
  basicstyle=\ttfamily\small,
  backgroundcolor=\color{codebg},
  frame=single,
  rulecolor=\color{codeframe},
  framerule=0.3pt,
  showstringspaces=false,
  keywordstyle=\color{codekey}\bfseries,
  commentstyle=\color{codecom}\itshape,
  breaklines=true,
  columns=fullflexible,
  captionpos=b,
  xleftmargin=0.5em,
  xrightmargin=0.5em,
}

% ---- Theorem environments ----
\newtheorem{definition}{Definition}
\newtheorem{theorem}{Theorem}
\newtheorem{lemma}{Lemma}
\newtheorem{remark}{Remark}

% ---- Title block ----
\title{Recursive Lattice zk-SNARKs for Instant Trustless Bootstrap:\\
       A 10-Millisecond Verification Path for Post-Quantum Blockchain Nodes\\
       \large{(v2 --- with zk-STARK--attested transparent setup)}}
\author{Quillon Graph Research\\
        \texttt{research@quillon.xyz}}
\date{May 13, 2026}

\begin{document}
\maketitle

\begin{abstract}
\noindent
We describe the design and phased implementation of a recursive zero-knowledge succinct
non-interactive argument of knowledge (zk-SNARK) for the Quillon Graph chain. The
construction collapses the validity of the entire chain tip into a single constant-size
proof verifiable in approximately five to ten milliseconds on commodity hardware.
A fresh node fetches a 5--150 KB proof from any peer, verifies it locally,
and immediately accepts the canonical state root --- enabling mining and transactions
within seconds rather than hours, without trusting a checkpoint signed by an operator.
The state commitment is a depth-256 sparse Merkle tree over the wallet--balance map;
the per-block transition is encoded as an R1CS circuit composing BLAKE3,
ML-DSA (Dilithium5) signature verification, a Merkle-path gadget, and a number-theoretic
transform (NTT)-based anchor election check; the recursion is realized via Nova
folding in the first deployment and is scoped to migrate to a lattice IVC scheme
(LatticeFold, LaBRADOR, or Greyhound) in 2028 so the entire trust chain becomes
post-quantum. The Nova structured reference string is generated by a deterministic
public algorithm and accompanied by a zero-knowledge STARK proof of correct generation,
removing the need for a multi-party trusted-setup ceremony and aligning with the
project's existing on-chain STARK infrastructure (4{,}469 LOC across
\texttt{crates/q-zk-stark/}, plus a 173~KB genesis-time encryption-SRS STARK proof
already on the live chain). We report the implementation status honestly: the
sparse Merkle tree is shipped in shadow mode; the gadget library (BLAKE3, Poseidon,
NTT, Dilithium5) totalling roughly 2{,}800 LOC is production-grade; the STARK
infrastructure (AIR framework, prover, verifier, batch prover, polynomial/FRI
commitments, blockchain-state circuit, wallet-privacy circuit) is approximately
70--80\% shipped; the $\delta$-circuit composition, Nova wrapper, lattice swap,
and wire protocol are scoped for Phases 1--4 with public mandatory verification
targeted for 2028.
\end{abstract}

\vspace{1em}
\noindent\rule{\linewidth}{0.4pt}
\tableofcontents
\noindent\rule{\linewidth}{0.4pt}
\vspace{0.5em}

% =========================================================================
\section{Introduction}
\label{sec:intro}

Joining a blockchain network for the first time is one of the worst experiences in
decentralized infrastructure. A new Quillon Graph node today must either
\textbf{(a)} replay the entire chain from genesis, which takes between five and
twelve hours depending on disk throughput and peer bandwidth, or \textbf{(b)} accept
a checkpoint snapshot signed by a bootstrap operator and trust that the signing key
has not been compromised. Option (a) is genuinely trustless but glacially slow.
Option (b) is fast but reintroduces a centralized trust assumption that the rest of
the protocol works hard to eliminate.

We argue that neither tradeoff is necessary. Recent advances in incrementally
verifiable computation (IVC)~\cite{Valiant2008,KothapalliSettyTzialla2021}, paired
with a state commitment carefully chosen to be both lattice-friendly and
hash-only (no elliptic-curve primitives buried in the consensus path),
make it possible to compress the validity of an entire chain into a single
constant-size proof that any node --- including a browser wallet --- can verify in
milliseconds.

The contribution of this whitepaper is the concrete design of such a proof for
Quillon Graph and a phased plan to deploy it without disrupting a live
multi-billion-dollar mainnet:

\begin{enumerate}[label=(\arabic*),leftmargin=*]
  \item A \textbf{depth-256 sparse Merkle tree} (SMT) over the wallet--balance map,
        using BLAKE3 leaf and node hashes with explicit domain separation. The SMT
        is implemented today in \texttt{crates/q-storage/src/balance\_smt.rs}
        ($\sim$600 LOC, twelve test cases, shipped in shadow mode).
  \item A \textbf{single-block state-transition circuit} ($\delta$-circuit) that
        composes the existing R1CS gadget library:
        BLAKE3 compression, ML-DSA (Dilithium5) signature
        verification, Poseidon, an NTT verifier, and a Merkle-path gadget over the SMT.
        Scoped for Phase~1; not yet implemented.
  \item A \textbf{Nova fold} of the $\delta$-circuit producing a constant-size,
        constant-verification-cost proof $\pi_n$ of the validity of the chain
        from genesis to height~$n$. Scoped for Phase~2.
  \item A \textbf{lattice migration path} (Phase~4, target Q1~2028) that
        replaces Nova on BN254 with LatticeFold~\cite{BonehChenTairi2024LatticeFold},
        LaBRADOR~\cite{BeullensSeiler2023LaBRADOR}, or
        Greyhound~\cite{NguyenSeiler2024Greyhound}, leaving the
        $\delta$-circuit unchanged and removing the trusted setup. The decision
        among the three candidates is deferred until Phase 2--3 benchmark data
        is available.
  \item A \textbf{progressive-archival UX} (\texttt{GET /api/v1/status/archive}
        and \texttt{X-QNK-Archive-*} response headers) that ships independently
        of the SNARK work, so the user-facing promise of an instantly useful
        node is honest about which queries are answered locally and which are
        still being backfilled.
\end{enumerate}

We are explicit throughout this document about what is implemented today and
what is forthcoming. The SMT is in production binaries in shadow mode (validating
parity with the legacy flat-hash \texttt{balance\_root\_v1} on every block).
The 2{,}792 LOC of gadgets in \texttt{crates/q-ivc/} cover BLAKE3 verification,
Poseidon, an NTT butterfly with the FIPS~204 negacyclic convention, and the full
Dilithium5 verification primitive (\texttt{compute\_az\_minus\_ct},
\texttt{high\_bits}, \texttt{use\_hint}, and the signed-norm helper). The
$\delta$-circuit composition, the Nova wrapper, the lattice swap, and the wire
protocol are aspirational; none are released yet. Promising what does not exist
is how cryptographic protocols quietly become unsafe in production.

The remainder of this paper is organized as follows. Section~\ref{sec:background}
reviews the necessary background: IVC, lattice succinct arguments, and the
specifics of the Quillon Graph chain. Section~\ref{sec:architecture} presents the
SMT and the $\delta$-circuit. Section~\ref{sec:performance} analyzes constraint
counts, prover throughput, verifier latency, and proof size.
Section~\ref{sec:lattice} discusses the lattice migration.
Section~\ref{sec:bootstrap} specifies the wire protocol and the progressive
archival UX. Sections~\ref{sec:security} and~\ref{sec:comparison} cover
security analysis and a comparison to related systems.
Section~\ref{sec:rollout} sets out the phased rollout.
Section~\ref{sec:status} states the present implementation status precisely.

% =========================================================================
\section{Background and Preliminaries}
\label{sec:background}

\subsection{The Quillon Graph chain}

Quillon Graph is a UTXO-style account chain whose consensus layer combines a
Narwhal-style DAG mempool~\cite{NarwhalTusk2022} with the
DAG-Knight anchor election~\cite{KaramanolakisDAGKnight2022}. Blocks are produced
at approximately one block per second on mainnet. The consensus layer chooses
an \emph{anchor block} per round using a verifiable number-theoretic transform
(NTT)-based randomness beacon. Transaction signatures use the ML-DSA scheme
standardized as FIPS~204~\cite{FIPS204} at parameter level 5 (``Dilithium5'').
Block headers are serialized in postcard
binary format and hashed with BLAKE3~\cite{BLAKE3spec}. The state commitment
in the current header is \texttt{balance\_root\_v1}, a flat BLAKE3 hash over the
sorted \texttt{(address, balance)} pairs:

\begin{lstlisting}[caption={Legacy state commitment in \texttt{crates/q-storage/src/lib.rs}.},label={lst:v1}]
let mut root_hasher = blake3::Hasher::new();
root_hasher.update(b"balance_root_v1");
for leaf in &leaf_hashes {     // O(N) in N wallets
    root_hasher.update(leaf);
}
*root_hasher.finalize().as_bytes()
\end{lstlisting}

This commitment is sound and cheap to maintain off-circuit, but proving membership
inside an R1CS circuit costs $\Theta(N)$ BLAKE3 compressions --- infeasible at
$N \approx 10^5$ wallets today and worse as the network grows.

\subsection{Incrementally verifiable computation}

Incrementally verifiable computation (IVC), introduced by Valiant in
2008~\cite{Valiant2008}, lets a prover maintain a single proof $\pi_n$ that
$F^n(z_0) = z_n$ for some step function $F$, updating it incrementally as
$F$ is applied. The verifier checks $\pi_n$ in time independent of $n$.
The original construction was theoretical; folding-scheme approaches such as
Nova~\cite{KothapalliSettyTzialla2021} and its refinement
HyperNova~\cite{KothapalliSetty2023HyperNova} have produced the first
implementations efficient enough for production-scale recursion.

Nova works over a pair of elliptic curves (typically BN254 and a cycle partner).
The folding step does not produce a SNARK at each step --- it produces a
\emph{relaxed} R1CS instance. A final SNARK is generated only at the moment
verification is required, by ``decompressing'' the relaxed instance through
a Spartan-style proof. This is what gives Nova its key property: the per-step
prover cost is linear in the step circuit, not in $n$, and the final verify cost
is independent of $n$.

For blockchain bootstrap, this is exactly the property we want.
The step function is ``apply one block to the state root.''
A 12-month-old chain with $\approx 31{,}500{,}000$ steps verifies in the same
time as a one-hour-old chain.

\subsection{Lattice succinct arguments}

Pairing-based SNARKs --- Groth16~\cite{Groth16},
Marlin~\cite{Marlin}, the Pickles construction in Mina~\cite{MinaPickles},
Halo~2~\cite{Halo2} --- rely on the hardness of the discrete log
or pairing problems on elliptic curves. These assumptions fall to a
sufficiently large quantum computer running Shor's algorithm~\cite{Shor1994}.
This is a problem for a chain whose transaction signatures
(Dilithium5) and key-encapsulation primitives (Kyber, FIPS~203~\cite{FIPS203}) are
already post-quantum: if the recursive proof itself rests on a quantum-broken
primitive, then the most critical piece of the trust chain --- the proof
that the entire history is valid --- has the worst quantum exposure in the
stack.

Lattice-based succinct arguments avoid this. LaBRADOR~\cite{BeullensSeiler2023LaBRADOR}
produces compact proofs for R1CS instances based on the Module Short Integer
Solution (Module-SIS) problem. LatticeFold~\cite{BonehChenTairi2024LatticeFold}
adapts the folding-scheme idea to lattices, providing IVC with Module-SIS
hardness. Greyhound~\cite{NguyenSeiler2024Greyhound} provides fast lattice-based
polynomial commitments and could be wrapped in a recursion harness.

All three are post-NIST PQC standardization era~\cite{NISTPQC2022} and
parameter-selectable to at least the NIST PQC level-1 security target.
The trade-off is implementation maturity: pairing-based recursive SNARKs
have production deployments going back years (Mina, Aleo); lattice IVC
deployments do not yet exist outside research code as of mid-2026.

\subsection{Why we will deploy Nova first and migrate later}

We choose to ship a pairing-based recursive SNARK (Nova on BN254) in
Phases~1--3, accepting a transient non-post-quantum dependency in the
recursive layer, and then migrate to lattice IVC in Phase~4. The reasoning is
operational rather than cryptographic:

\begin{itemize}[leftmargin=*]
  \item The $\delta$-circuit is the same in both flavors. The proof
        system swap is a localized change of about 800 LOC in
        \texttt{crates/q-ivc/src/recursion/}; everything else --- the
        SMT, the gadgets, the API surface, the bootstrap flow --- stays put.
  \item Debugging a 440-million-constraint composition circuit
        \emph{and} a brand-new lattice IVC implementation simultaneously
        is a recipe for soundness errata. Production-mature Nova lets us
        validate the circuit first.
  \item The recursive proof during Phases~1--3 is advisory --- blocks
        are still validated block-by-block on every node. A quantum
        attacker forging a Nova proof during this window would not
        compromise the chain; the proof simply would not be accepted
        when re-verified against block-level rules.
  \item By the Phase 4 target (Q1 2028), at least one of the three
        lattice candidates will have a vetted reference implementation
        suitable for production. We pick based on benchmark data at
        that time, not now.
\end{itemize}

% =========================================================================
\section{System Architecture}
\label{sec:architecture}

\subsection{State commitment: \texttt{balance\_root\_v2} (sparse Merkle tree)}

The first structural change is the migration from the flat
\texttt{balance\_root\_v1} commitment to a sparse Merkle tree
(\texttt{balance\_root\_v2}). The SMT is keyed by the 32-byte BLAKE3 of the
public key (the wallet address); leaves carry the 16-byte little-endian
\texttt{u128} balance. The tree has depth 256, matching the address length.
Internal nodes use BLAKE3 with explicit domain separation:

\[
  \begin{aligned}
    \mathrm{leaf}(a, b) &= \mathrm{BLAKE3}(\mathtt{"smt\_leaf\_v2"} \,\|\, a \,\|\, b_{\mathrm{LE}}) \\
    \mathrm{node}(L, R) &= \mathrm{BLAKE3}(\mathtt{"smt\_node\_v2"} \,\|\, L \,\|\, R)
  \end{aligned}
\]

The empty-subtree hashes
$E_d = \mathrm{node}(E_{d+1}, E_{d+1})$ are precomputed for all depths
$d = 0, \dots, 256$, anchored by $E_{256} = \mathrm{leaf}(\mathbf{0}, 0)$.
When a sibling along a path equals $E_d$ the storage layer omits it from
the database; on the wire it is signaled by a single bit in a 256-bit
empty-subtree bitmap. This is the standard ``Compact Sparse Merkle Tree''
construction in proof-friendly form.

\textbf{Implementation status:} \texttt{crates/q-storage/src/balance\_smt.rs}
is approximately 743 LOC and is shipped in shadow mode in
v10.9.x: every block updates both \texttt{balance\_root\_v1} and the
SMT in the same RocksDB write batch and asserts identity between the
two roots (so any divergence is immediately visible in mainnet logs).
The mandatory switch from v1 to v2 in the block header is gated on a
future activation height \texttt{BALANCE\_ROOT\_V2\_HEIGHT}, which is
not yet set.

\begin{lstlisting}[caption={Public API of the shipped sparse Merkle tree, \texttt{crates/q-storage/src/balance\_smt.rs}.},label={lst:smtapi}]
pub struct BalanceSmt { /* ... */ }

impl BalanceSmt {
    pub fn new(db: Arc<DB>) -> Result<Self>;
    pub async fn rebuild_from_balances(
        &self, balances: &HashMap<[u8;32], u128>
    ) -> Result<[u8;32]>;
    pub async fn update(
        &self, addr: &[u8;32], balance: u128
    ) -> Result<[u8;32]>;
    pub async fn batch_update(
        &self, updates: &[([u8;32], u128)]
    ) -> Result<[u8;32]>;
    pub async fn prove(&self, addr: &[u8;32]) -> Result<SmtProof>;
    pub fn root(&self) -> [u8;32];
}

pub struct SmtProof {
    pub addr: [u8;32],
    pub balance: u128,
    pub siblings: [[u8;32]; 256],
    pub empty_bitmap: [u8; 32],  // bit i = 1 iff siblings[i] == E_i
}

impl SmtProof {
    pub fn verify(&self, expected_root: &[u8;32]) -> bool;
}
\end{lstlisting}

\subsection{The $\delta$-circuit: one R1CS proof per block transition}

The $\delta$-circuit is the heart of the construction. It takes as
\emph{public inputs} the previous and next state roots, the block header
hash, and the block height. It takes as \emph{private witnesses} the
entire block body, the Merkle paths into both state roots for every
touched wallet, the Dilithium5 signature verification witnesses, and
the NTT anchor witness. The circuit is satisfied if and only if the
block validly transitions \texttt{state\_root\_prev}~$\to$~\texttt{state\_root\_next}.

\begin{definition}[$\delta$-circuit]
\label{def:deltacircuit}
Let $s_{n}$ denote the SMT root at height $n$ and $B_{n+1}$ the block
produced at height $n+1$. The $\delta$-circuit is an R1CS predicate
\[
  \delta\big(s_n,\, B_{n+1},\, s_{n+1}\big) \in \{0,1\}
\]
satisfied if and only if $B_{n+1}$ is a valid block under all
consensus rules and applying its transactions, coinbase, and emission
to the wallet--balance map committed by $s_n$ yields the map committed
by $s_{n+1}$.
\end{definition}

The predicate composes the existing gadgets:

\begin{enumerate}[leftmargin=*]
  \item $h_{\mathrm{hdr}} = \mathrm{BLAKE3}(\mathrm{header\_bytes})$,
        enforced via \texttt{Blake3Gadget::verify\_hash}.
  \item For each transaction $\mathrm{tx} = (\mathrm{from}, \mathrm{to}, a, f, n, \sigma, \mathrm{pk})$:
        \begin{enumerate}[label*=(\alph*),leftmargin=*]
          \item $\mathrm{Verify}_{\mathrm{Dilithium5}}(\mathrm{pk}, \mathrm{msg}(\mathrm{tx}), \sigma)$,
                via \texttt{DilithiumVerifierGadget::verify}.
          \item Membership of $(\mathrm{from}, b_{\mathrm{from}})$ in $s_n$, via
                \texttt{MerklePathGadget::enforce\_membership}.
          \item Membership of $(\mathrm{to}, b_{\mathrm{to}})$ in $s_n$, same gadget.
          \item Balance sufficiency $b_{\mathrm{from}} \ge a + f$
                (range check on $b_{\mathrm{from}} - a - f$ via \texttt{enforce\_norm\_bound}).
          \item Updated $(\mathrm{from}, b_{\mathrm{from}} - a - f)$ membership in
                an intermediate root.
          \item Updated $(\mathrm{to}, b_{\mathrm{to}} + a)$ membership in the
                root that follows the from-side update.
        \end{enumerate}
  \item Coinbase emission satisfies the era-step schedule
        $e \le R(\mathrm{height})$, where $R$ is implemented as a piecewise
        lookup over the four-year halving boundaries.
  \item NTT-based anchor election validates the block producer's claim
        of being the anchor for this round, via \texttt{NttVerifierGadget}.
  \item After all transactions and the coinbase have been applied
        sequentially, the resulting current root must equal the public
        \texttt{state\_root\_next}.
\end{enumerate}

\begin{remark}
\label{rem:perblockfreshness}
The $\delta$-circuit is large but not enormous. We accept a high per-block
prover cost (Section~\ref{sec:performance}) because the prover lives
on dedicated infrastructure and runs once per block; the verifier runs
on every wallet and every new node, and we want it fast.
\end{remark}

\subsection{Recursion: Nova folding}

With $\delta$ defined as a step circuit, the recursive proof is
constructed by Nova's fold:
\[
  \pi_{n+1} = \mathrm{Nova.Fold}\big(\delta,\ s_n,\ B_{n+1},\ \pi_n\big),
  \qquad \pi_0 = \mathbf{1}.
\]
Each fold consumes the previous proof and the new block's witness and
produces a fresh proof of size and verification cost independent of $n$.
Genesis nodes (Epsilon, Beta, Gamma, Delta) maintain
$(s_n, \pi_n)$ alongside the block database and serve
$(s_n, \pi_n, \mathrm{header}_n)$ over the wire protocol described in
Section~\ref{sec:bootstrap}.

A new node executes the algorithm in Algorithm~\ref{alg:bootstrap}.

\begin{algorithm}[t]
\caption{Proof-bootstrap of a fresh node}
\label{alg:bootstrap}
\begin{algorithmic}[1]
\Require Peer URL $p$, expected genesis state root $s_0$
\State $(s_t, \pi_t, H_t) \gets \mathrm{HTTP\_GET}(p, \mathtt{/api/v1/proof/tip})$
\State $\mathrm{valid} \gets \mathrm{VerifyRecursive}(\pi_t, s_t, H_t.\mathrm{height})$
\If{$\neg\mathrm{valid}$}
  \State \textbf{retry} with a different peer
\EndIf
\State \textbf{assert} $H_t.\mathrm{state\_root} = s_t$
\State Accept $s_t$ as canonical; enable mining and transaction acceptance
\State Begin archive backfill in background (not on critical path)
\end{algorithmic}
\end{algorithm}

Steps 1--7 complete in well under one second; the verify call (line 2)
is the binding constraint at approximately 5--10 ms.

% =========================================================================
\section{Transparent Setup via zk-STARK Attestation}
\label{sec:stark-setup}

\subsection{The trusted-setup problem}

Nova and Groth16 both consume a structured reference string (SRS) at proving
and verification time. The SRS is generated by sampling a secret trapdoor $\tau$
and emitting public group elements derived from $\tau$. Anyone who knows $\tau$
can forge proofs. The standard mitigation is a multi-party computation (MPC)
ceremony: $N$ participants each contribute randomness, mix it into the SRS,
and discard their share. Soundness requires that at least one participant
genuinely discarded their share. The Powers of Tau ceremonies behind
\cite{Groth16} and most Zcash deployments adopted this protocol with hundreds
of participants. The ceremony is correct in principle but operationally
unattractive: it is slow to coordinate, must be repeated for every new
circuit, and requires real-world trust in the participant pool.

\subsection{The STARK-attested alternative}

We avoid the ceremony entirely. The SRS is generated by a deterministic
public algorithm $\mathsf{gen}: \mathsf{seed} \mapsto \mathsf{SRS}$, where
$\mathsf{seed}$ is a public unpredictable quantity (in our case, a future
block hash determined after the SRS-generation transaction is mined --- a
RANDAO-style construction). A zero-knowledge STARK proof
$\pi_{\mathrm{setup}}$ attests that $\mathsf{SRS}$ was produced by running
$\mathsf{gen}$ on $\mathsf{seed}$. Any validator can re-run $\mathsf{gen}$
locally and verify $\pi_{\mathrm{setup}}$; if both check, the SRS is accepted.

The trust assumption collapses to:
\begin{enumerate}[leftmargin=2em,topsep=0pt]
\item \textbf{STARK soundness} (random-oracle model on BLAKE3
      \cite{BLAKE3spec,BSBHR2018STARK}), and
\item \textbf{Seed unpredictability} (the seed is committed before any party
      who could influence it learns its value).
\end{enumerate}
Both assumptions are standard, well-studied, and require zero coordination
between human participants.

\vspace{0.5em}
\begin{verbatim}
   public seed s ──┬──> gen(s)  ──────> SRS
                   │
                   └──> STARK_Prove(gen, s, SRS) ──> π_setup
                                                     │
   validator: re-run gen(s), check SRS matches,      ▼
              verify π_setup with public verifier  trust SRS
\end{verbatim}
\vspace{0.5em}

\subsection{Existing infrastructure}

We are not adding zk-STARK to the project for this work. The pattern is
already in production:

\begin{itemize}[leftmargin=2em,topsep=0pt,itemsep=0pt]
\item \texttt{crates/q-zk-stark/} (4{,}469 lines of Rust) ships:
  \texttt{air.rs} (Algebraic Intermediate Representation, 429 LOC),
  \texttt{stark\_prover.rs} (442 LOC), \texttt{stark\_verifier.rs}
  (548 LOC), \texttt{batch\_prover.rs} (662 LOC, batched STARK proofs over
  multiple statements), \texttt{polynomials.rs} (549 LOC, FRI-style polynomial
  commitments), \texttt{blockchain\_state\_circuit.rs} (410 LOC, a STARK
  circuit over blockchain state transitions), \texttt{wallet\_privacy\_stark.rs}
  (530 LOC, deployed wallet-privacy proofs), and a \texttt{gpu/} subtree for
  hardware acceleration.
\item \texttt{crates/q-storage/src/encryption\_zkstark.rs} (562 LOC) exposes
  \texttt{EncryptionKeyProof} --- a STARK proof binding the encryption key
  schedule to a public seed. The proof is generated once per chain at
  genesis, validated by every participating node.
\item The artifact \texttt{data-mainnet-genesis/encryption.zkstark}
  (176{,}832 bytes) is a production STARK proof currently committed to
  the chain. The transparent-setup pattern has been live since
  \texttt{v1.0.43-beta}.
\end{itemize}

The Nova-SRS attestation reuses the same AIR framework and the same prover/verifier
pair. The circuit that the STARK proves is the Nova-SRS generator $\mathsf{gen}$
itself: a deterministic sequence of field-multiplications and group exponentiations.
No new cryptographic constructions are required --- only a new AIR description
of $\mathsf{gen}$ and a one-time prover invocation.

\subsection{What this changes operationally}

The Phase~2 deployment (Section~\ref{sec:phased}) ships the Nova-attested SRS from
day one. There is no separate MPC-ceremony phase. The on-chain artifact
\texttt{data-mainnet-genesis/nova-srs.zkstark} (analogous to the existing
\texttt{encryption.zkstark}) is generated by the Quillon Graph maintainers,
verified by every validator and every fresh-bootstrap node, and remains
auditable indefinitely. If a flaw is later found in the chosen lattice scheme
or in Nova itself, the STARK-attested SRS for the next generation can be
re-generated and re-attested without re-running a coordinated MPC ceremony.

\subsection{What this does \emph{not} change}

The STARK attestation removes the ceremony but does not make the underlying
proof system post-quantum. Nova on BN254 still depends on the elliptic-curve
discrete logarithm assumption, which Shor's algorithm \cite{Shor1994} can
break given a sufficiently capable quantum computer. Phase~4
(Section~\ref{sec:lattice-migration}) addresses post-quantum by swapping the
underlying scheme; the STARK-attested setup pattern carries over unchanged
to any lattice scheme that requires public parameters.

% =========================================================================
\section{Performance Analysis}
\label{sec:performance}

\subsection{Constraint count breakdown}

The cost of a single $\delta$-circuit step dominates the prover budget.
Table~\ref{tab:constraints} breaks down the constraint count for a
block carrying $K = 100$ transactions and one coinbase, in line with
current mainnet block content.

\begin{table}[h]
\centering
\begin{tabular}{lrr}
\toprule
\textbf{Component} & \textbf{Per-instance} & \textbf{Block total ($K=100$)} \\
\midrule
Block-header BLAKE3 hash               & 50{,}000        & 50{,}000 \\
Per-tx Dilithium5 signature verify     & 1{,}500{,}000   & 150{,}000{,}000 \\
Per-tx Merkle paths (4 $\times$ depth-256, BLAKE3) & 2{,}360{,}000 & 236{,}000{,}000 \\
Per-tx balance arithmetic + range      & 10{,}000        & 1{,}000{,}000 \\
Coinbase Merkle path + emission lookup & 5{,}000{,}000   & 5{,}000{,}000 \\
NTT anchor verification                & 50{,}000{,}000  & 50{,}000{,}000 \\
\midrule
\textbf{Total}                         &                 & \textbf{$\approx$ 442{,}050{,}000} \\
\bottomrule
\end{tabular}
\caption{R1CS constraint count breakdown for one $\delta$-circuit step, $K = 100$ transactions.
The dominant terms are Dilithium5 signature verification and the four Merkle paths per
transaction. Numbers are upper-bound estimates derived from the existing
\texttt{Blake3Gadget} per-compression cost (approximately 2{,}300 constraints)
and from published Dilithium R1CS analyses.}
\label{tab:constraints}
\end{table}

Two notes on the table. First, the four Merkle paths per transaction
($\mathrm{from\_prev}$, $\mathrm{to\_prev}$, $\mathrm{from\_next}$,
$\mathrm{to\_next}$) at 256~$\times$ BLAKE3-compression each costs about
590{,}000 constraints per path. Second, the BLAKE3 cost dominates the
Merkle gadget; replacing it with a SNARK-friendly hash such as
Poseidon~\cite{Poseidon} would compress this column by roughly two orders
of magnitude but would also make the SMT root no longer reproducible by a
BLAKE3-only light client. We have chosen to stay BLAKE3 for v2 and to
revisit a Poseidon-rooted ``v3'' state commitment in 2027 if Phase 2--3
benchmarks demand it.

\subsection{Prover throughput target}

To keep up with mainnet's one-block-per-second cadence, the fold step
must complete within one wall-clock second per block on the genesis
prover hardware. Nova's per-step prover cost scales roughly linearly
with the step circuit's constraint count; a 442~M constraint step on
a 48-core Epsilon-class machine is on the edge of feasibility. We
have three optimization levers in reserve:

\begin{enumerate}[leftmargin=*]
  \item \textbf{Batched signature verification.} Verify $b$ signatures
        in a batched gadget rather than $b$ independent gadget invocations,
        reusing the same Number-Theoretic Transform tables. Saves on the
        order of 40\% per signature for $b = 4$.
  \item \textbf{Multi-block folds.} Fold $m$ blocks into one Nova step
        (still valid; Nova does not require step granularity to match
        block granularity). This amortizes setup overhead.
  \item \textbf{Distributed proving.} Phase~3 onwards, distribute the
        per-block fold across multiple validator nodes so no single
        machine has to keep up with mainnet alone.
\end{enumerate}

If none of these are sufficient, we fall back to producing $\pi$ every
$m$ blocks rather than every block, with $m \le 10$. A bootstrap proof
that lags the tip by ten seconds is still acceptable user experience.

\subsection{Verifier latency}

Verifier cost is the metric that ultimately matters for the
user-facing promise. Table~\ref{tab:verifylatency} summarizes the
targets across deployment surfaces.

\begin{table}[h]
\centering
\begin{tabular}{lrr}
\toprule
\textbf{Surface} & \textbf{Target latency} & \textbf{Notes} \\
\midrule
Desktop (Intel/AMD x86-64, 8 cores)  & $\le$ 5 ms   & Native arkworks verifier \\
Laptop (M-series ARM)                & $\le$ 10 ms  & Native arkworks verifier \\
Browser WASM (mobile or desktop)     & $\le$ 250 ms & wasm-bindgen + IndexedDB cache \\
Embedded (Raspberry Pi 5)            & $\le$ 50 ms  & Native arkworks verifier \\
\bottomrule
\end{tabular}
\caption{Verifier latency targets across deployment surfaces. Targets
are conservative; published Nova-on-BN254 benchmarks consistently
report sub-5-ms verify times on commodity desktop hardware.}
\label{tab:verifylatency}
\end{table}

The advertised ``10-millisecond verification'' in the title of this
paper refers to the laptop target. The desktop number is faster; the
browser number is slower but still well within the latency budget of
a wallet's first-load operation.

\subsection{Proof size}

Nova proofs over BN254 are approximately 5 KB when finalized through a
Spartan compression step. Lattice proofs are larger: LaBRADOR's
published figures put R1CS proofs in the 50--150 KB range depending on
the constraint count. Both fit comfortably in a single HTTP response
and well below the gossipsub message-size limits we have configured.

% =========================================================================
\section{The Lattice Migration}
\label{sec:lattice}

\subsection{Why migrate}

Once Phases~1--3 are stable, the only non-post-quantum primitive remaining
in the consensus path is the elliptic-curve folding scheme in Nova.
This is acceptable transient state but not acceptable long-term. The
2028 lattice migration is what closes that gap.

\subsection{Setup options across the Phase~2--4 trajectory}

Combining the Nova/BN254 choice (Phases~1--3) with the STARK-attested setup
(Section~\ref{sec:stark-setup}) and the candidate lattice schemes for Phase~4
gives four operationally distinct configurations. Table~\ref{tab:setup-options}
compares them. The middle two rows describe what we actually ship in
Phases~2 and 3; the bottom two are the post-quantum migration targets.

\begin{table}[h]
\centering
\small
\begin{tabularx}{\linewidth}{l X X X X}
\toprule
\textbf{Scheme} & \textbf{Verify} & \textbf{Proof size} & \textbf{Setup} & \textbf{PQ-safe} \\
\midrule
Nova/BN254 (MPC ceremony)         & $\sim$5 ms                  & $\sim$5 KB              & MPC, $N$ participants                       & No \\
\textbf{Nova/BN254 + STARK-attested SRS} (Phase~2--3)
                                  & $\sim$5 ms + STARK verify   & $\sim$5 KB + $\sim$50 KB STARK
                                                                                          & Transparent (public algorithm + STARK)
                                                                                                                                        & No \\
LatticeFold (Phase~4 candidate A) & $\sim$50 ms (target 10 ms)  & 50--100 KB              & None                                         & Yes \\
LaBRADOR-IVC + wrapper            & $\sim$30 ms (target 10 ms)  & 50--150 KB              & None                                         & Yes \\
Greyhound (poly commit + adapt.)  & $\sim$30 ms (target 10 ms)  & 100 KB+                 & None                                         & Yes \\
\bottomrule
\end{tabularx}
\caption{Setup and verification trade-offs across deployment phases.
The boldface row is what Phase~2--3 ship. The bottom three rows are the
Phase~4 candidates, all of which avoid setup entirely by virtue of their
algebraic construction (no SRS, no trapdoor). The decision among the
three lattice candidates is deferred until Phase~2--3 benchmark data
is collected.}
\label{tab:setup-options}
\end{table}

The Nova-on-BN254 path is not post-quantum, but the STARK attestation of
its setup is. Once Phase~4 swaps to a lattice scheme, the entire chain
becomes post-quantum end-to-end, and the STARK attestation is no longer
needed for that proof system (the lattice schemes require no setup at all).
The STARK pattern itself remains useful for other places in the stack
(Section~\ref{sec:stark-elsewhere}).

\subsection{What stays the same across the swap}

The decision to migrate is meaningful only because the swap surface is
narrow. The $\delta$-circuit definition in
\texttt{crates/q-ivc/src/circuits/delta\_block.rs} is unchanged.
\texttt{balance\_root\_v2} is unchanged. The wire protocol shape
(Section~\ref{sec:bootstrap}) is unchanged --- only the
\texttt{proof\_version} string and the byte encoding of \texttt{proof\_b64}
differ. The mandatory verification block-height gate logic is unchanged.
The advisory-vs-mandatory rollout sequence is unchanged.

What changes is approximately 800 LOC in
\texttt{crates/q-ivc/src/recursion/}: the \texttt{StepCircuit}
implementation is re-emitted against the lattice scheme's API; the
folder orchestrator is re-implemented; the verifier is re-implemented;
the public parameters file is regenerated.

\subsection{Decision pathway}

We commit explicitly to the following sequence:

\begin{enumerate}[leftmargin=*]
  \item \textbf{Late 2026:} prototype each of the three candidates as a
        non-IVC lattice argument of single-block validity. Measure
        prover time, verifier time, and proof size.
  \item \textbf{Q1 2027:} decide. Publish a short technical report
        committing to one of the three.
  \item \textbf{Q1--Q3 2027:} build the lattice IVC harness around the
        chosen scheme.
  \item \textbf{Q4 2027:} side-by-side testnet running Nova and lattice
        in parallel; identical $\delta$-circuit; compare proof times
        and verifier times under real block flow.
  \item \textbf{Q1 2028:} production swap.
\end{enumerate}

% =========================================================================
\section{Beyond Setup: zk-STARK Throughout the Stack}
\label{sec:stark-elsewhere}

The STARK-attested SRS (Section~\ref{sec:stark-setup}) is the first place we
apply zk-STARK in the recursive-bootstrap stack, but it is not the last.
The same \texttt{crates/q-zk-stark/} infrastructure composes naturally with
several other parts of the architecture. We list them here in order of
deployment plausibility and label each as \emph{shipped today},
\emph{near-term work}, or \emph{future research}.

\subsection{Wallet privacy proofs --- shipped today}

\texttt{crates/q-zk-stark/src/wallet\_privacy\_stark.rs} (530 LOC) is
production code today. A wallet can prove ``I own at least $X$ QUG''
without revealing the exact balance. The proof is a STARK over a
range-membership AIR, verifiable in pure native Rust on any node and in
WebAssembly via the same bindings as the recursive SNARK verifier
(Section~\ref{sec:bootstrap}). This existing capability is independent of
the SRS-attestation pattern but uses the same underlying AIR framework.
The relevance to recursive SNARK: it demonstrates that mature, deployed
STARK verification in Quillon Graph is not theoretical --- nodes already
verify STARKs on the consensus path.

\subsection{Batched Dilithium signature verification --- near-term work}

The dominant per-block cost in the $\delta$-circuit is signature
verification. Each Dilithium5 signature requires approximately 1.5~M
R1CS constraints; a 100-transaction block accumulates 150~M constraints
from signatures alone (Section~\ref{sec:performance}). A natural
optimization: replace in-circuit signature verification with a
\emph{batched STARK proof} that asserts ``these $K$ signatures all
verify correctly,'' consumed by the $\delta$-circuit as a single
attestation. The reduction is dramatic: STARK signature verification
scales nearly linearly in $K$ outside the circuit (a few hundred
microseconds per signature on dedicated hardware) and the $\delta$-circuit
swallows the result as a constant-size constraint.

The infrastructure is already partly in place:
\texttt{crates/q-zk-stark/src/batch\_prover.rs} (662 LOC) implements
batched STARK proving over multiple statements with shared trace
columns. We expect this to reduce $\delta$-circuit constraint count
by 30--50\%, bringing the Nova prover budget per block from $\sim$60~s
toward the target of 1~s with commodity hardware.

\subsection{Light-client mobile path --- near-term work}

The recursive SNARK proof verifier targets $\leq 250$~ms in WebAssembly
(Section~\ref{sec:performance}). For extremely resource-constrained
clients --- low-end mobile devices, embedded peers, IoT clients ---
even 250~ms can be too much. A STARK proof of ``the most recent $N$
block headers are valid'' is verifiable in pure JavaScript (no WASM)
in tens of milliseconds, at the cost of trusting only the last $N$
blocks rather than the entire chain history. This is strictly weaker
than the recursive SNARK guarantee but acceptable for a light wallet
that only needs to know the current state, not the historical chain.

The pattern parallels \texttt{wallet\_privacy\_stark.rs}: a small,
mobile-friendly STARK proof, generated by genesis nodes, served alongside
the recursive proof for clients that opt into the lighter mode.

\subsection{Cross-chain bridge attestations --- future work}

Recursive SNARK proofs over BN254 are awkward to verify on chains that
use different elliptic curves. Verifying a Nova proof on Ethereum, for
example, requires implementing BN254 pairing arithmetic in Solidity ---
expensive in gas and a notable attack surface. STARK proofs use only
hash functions and are agnostic to the destination chain's curve
choice. A STARK proof of ``the Quillon Graph state root at height $H$
is $R$'' can be verified on Ethereum (via the Cairo-style verifier
pattern), Bitcoin (via Tapscript with appropriate hash gadgets), or any
chain that supports hash evaluation in its scripting language.

This is the natural way to bridge Quillon Graph state to other chains
without trusting a multi-sig committee: the destination chain verifies
a STARK over the bridged state root and accepts it as canonical. We
anticipate this becoming relevant in 2027 when bridging questions arise
in production. \texttt{crates/q-zk-stark/src/blockchain\_state\_circuit.rs}
(410 LOC) is the existing STARK circuit over blockchain state transitions
that would underpin this.

\subsection{Activation-height attestations --- future work}

When a consensus rule activates at a known block height (e.g.,
\texttt{BALANCE\_ROOT\_V2\_HEIGHT}), a STARK proof attesting ``the
state transition at this height correctly applies the rule activation''
provides an independent audit trail. Adversarial activations (a
malicious miner producing a block at the activation boundary with
inconsistent semantics) are detectable by any client that verifies
the activation STARK. The pattern reuses
\texttt{blockchain\_state\_circuit.rs}.

\subsection{Encryption SRS proof --- shipped today, foundational reference}

We close this section by emphasizing that the transparent-setup pattern
of Section~\ref{sec:stark-setup} is not new infrastructure. The
encryption SRS at chain genesis was produced by precisely this pattern
and lives in production today as a 176{,}832-byte STARK proof at
\texttt{data-mainnet-genesis/encryption.zkstark}. Every full node
verifies it at startup. The Nova-SRS attestation is the same engineering
applied to a different generator circuit.

% =========================================================================
\section{Bootstrap Protocol and Progressive Archival}
\label{sec:bootstrap}

\subsection{Wire protocol}

The bootstrap surface is three additions to the existing HTTP and
gossipsub layers.

\paragraph{\texttt{GET /api/v1/proof/tip}.}
Returns the latest folded proof and the block header it commits to.
No authentication --- the proof is publicly verifiable, so the API
exposes no authority.

\begin{lstlisting}[language=json,caption={\texttt{GET /api/v1/proof/tip} response body.},label={lst:proof-tip-resp}]
{
  "tip_height": 11400000,
  "state_root": "0xabcd...",
  "block_header": {
    "height": 11400000,
    "parent_hash": "0x...",
    "tx_root": "0x...",
    "state_root": "0xabcd...",
    "timestamp": 1771761600,
    "producer_id": 123,
    "anchor_validator": "..."
  },
  "proof_version": "nova-bn254-v1",
  "proof_size_bytes": 51232,
  "proof_b64": "<base64-encoded proof bytes>"
}
\end{lstlisting}

\paragraph{\texttt{POST /api/v1/proof/verify}.}
Optional helper for clients without an arkworks verifier compiled in
(notably browser wallets that prefer not to bundle WASM crypto).
Submits a proof and returns a boolean plus a verify-time measurement.
This endpoint is a convenience; the trustless path is for the client
to verify locally.

\paragraph{Gossipsub topic \texttt{/qnk/mainnet-genesis/recursive-proof}.}
Genesis nodes publish $(s_n, \pi_n, \mathrm{header}_n)$ on a new
gossipsub topic with one message per block. Peers maintain ``latest
seen proof'' and forward only higher-height proofs. Rate-limited to
one message per peer per block.

\subsection{Bootstrap-mode CLI flag}

The Quillon Graph node binary gains a single new flag:

\begin{lstlisting}
--bootstrap-mode=proof       # Download proof, verify, start immediately
--bootstrap-mode=checkpoint  # Existing behavior: trust a signed snapshot
--bootstrap-mode=genesis     # Existing behavior: full replay from height 1
\end{lstlisting}

\texttt{proof} is opt-in throughout Phase~3 (the advisory window) and
becomes the default in Phase~5 (mandatory verification).

\subsection{Progressive archival UX}

The proof-bootstrap delivers a verified state root, but it does not
deliver historical blocks. An archive backfill runs in parallel after
bootstrap. Historical-block queries against blocks not yet
backfilled need a graceful UX.

We add a status endpoint and three response headers:

\begin{lstlisting}[language=json,caption={\texttt{GET /api/v1/status/archive} response body.},label={lst:status-archive}]
{
  "tip_height": 11400000,
  "lowest_indexed_height": 4321001,
  "archive_complete": false,
  "archive_progress_pct": 37.9,
  "archive_eta_seconds": 64800,
  "blocks_per_sec_recent": 175,
  "verified_proof_height": null
}
\end{lstlisting}

Every historical-lookup endpoint (\texttt{/blocks/{h}}, \texttt{/tx/{id}},
\texttt{/receipts/{h}}) carries three response headers:
\texttt{X-QNK-Archive-Status: complete|backfilling},
\texttt{X-QNK-Archive-Progress: <indexed>/<tip>}, and
\texttt{X-QNK-Archive-Lowest-Indexed: <height>}. Queries for heights
below \texttt{lowest\_indexed\_height} return HTTP~202 (Accepted) with
a body describing where the data will be available, instead of HTTP~404.
Wallet and explorer UIs render this as ``Block X not yet indexed,
ETA T'' with a clearly marked opt-in proxy to a fully-indexed peer.

This UX ships in v10.9.16 \emph{independently of the SNARK work}:
the existing two-phase backfill already produces
\texttt{lowest\_indexed\_height}, so the endpoint and headers light up
immediately and are useful even before the proof bootstrap exists.

% =========================================================================
\section{Security Analysis}
\label{sec:security}

\subsection{Soundness}

The soundness of the bootstrap argument decomposes into three layers.

\textbf{Layer 1: the underlying proof system.} Nova on BN254 is sound
under the discrete-logarithm assumption on the BN254 curve and the
Random Oracle Model. We will rely on the published Nova literature
through Phase~3 and on the published lattice scheme analyses in
Phase~4. We commit to not modifying the proof system internals beyond
the published reference implementation.

\textbf{Layer 2: the $\delta$-circuit's correctness.} The circuit is
sound if and only if it faithfully encodes the block-validity predicate.
This is the layer where bugs can be introduced and is the principal
soundness risk. Mitigation: an adversarial-witness test suite covering
all five well-defined fail modes (invalid signature, insufficient
balance, state-root mismatch, replay-nonce reuse, over-emission
coinbase) is mandatory before mandatory verification at activation.

\textbf{Layer 3: the wire path.} The new node fetches a proof from
an arbitrary peer. The peer can return whatever bytes it likes. A
forged proof is rejected by the local verifier; a valid proof for a
different state root than the claimed one is rejected by the
post-verification consistency check at line~6 of
Algorithm~\ref{alg:bootstrap}. A network-level adversary who controls
all peers can stall the bootstrap but cannot induce the new node to
accept a wrong state root.

\subsection{Setup trust assumptions}
\label{sec:setup-trust}

Section~\ref{sec:stark-setup} described the transparent-setup pattern in
constructive detail. Here we make its trust assumptions precise and contrast
them with the alternatives:

\begin{itemize}[leftmargin=2em,itemsep=0pt]
\item \textbf{Naive Nova (rejected).} A single dealer generates the SRS,
      keeping the trapdoor $\tau$. Trust assumption: ``the dealer discarded
      $\tau$.'' Operationally untenable for a $\$2$ billion mainnet.
\item \textbf{Nova + MPC ceremony (industry standard, rejected here).}
      $N$ contributors mix randomness. Trust assumption: ``at least one
      of $N$ participants is honest \emph{and} actually discarded their
      contribution.'' Logistically expensive; requires coordinated public
      ceremony; must be re-run if the circuit changes.
\item \textbf{Nova + STARK-attested transparent setup (what we ship).}
      Trust assumption: ``STARK soundness in the Random Oracle Model on
      BLAKE3 \cite{BLAKE3spec,BSBHR2018STARK} \emph{and} the public seed
      is unpredictable at SRS-generation time.'' No coordination required;
      the proof is reproducible; mid-chain re-attestation is a single
      transaction.
\item \textbf{LatticeFold / LaBRADOR / Greyhound (Phase~4).} No setup
      at all. Trust assumption reduces to the chosen scheme's hardness
      assumption (Module-SIS / RLWE).
\end{itemize}

The STARK assumption stack is non-circular: STARK soundness only needs
BLAKE3 to behave as a random oracle. We already rely on BLAKE3 throughout
the chain for state roots, block hashes, and SMT leaves; the STARK
attestation does not introduce a new cryptographic primitive into the
trust model.

The STARK-attested pattern is already in production for a related setup:
the genesis encryption key schedule (\texttt{data-mainnet-genesis/encryption.zkstark},
176{,}832 bytes, generated by \texttt{crates/q-storage/src/encryption\_zkstark.rs}).
The Nova-SRS attestation is the same pattern applied to a different generator
circuit, and reuses every component of \texttt{crates/q-zk-stark/}.

\subsection{Phased deployment with advisory window}

The single most important mitigation against $\delta$-circuit bugs is
operational: the proof is \textbf{advisory} for at least six months on
mainnet before any node makes a consensus-affecting decision based on
the proof's validity. During the advisory window:

\begin{itemize}[leftmargin=*]
  \item Every block is still validated block-by-block on every node.
  \item The proof is published and verified, but its acceptance or
        rejection does not influence consensus.
  \item Mismatches between block-by-block validation and proof
        verification are logged as alerts; any non-zero rate of
        mismatches \emph{must} block the transition to mandatory
        verification.
\end{itemize}

Only after the advisory window completes with zero soundness
discrepancies does the chain activate mandatory verification at a
publicly announced block height.

\subsection{Lattice parameter selection}

When Phase~4 lands, parameters are chosen to meet at least NIST PQC
level~1 security against both classical and quantum adversaries, per
\cite{NISTPQC2022}. Specific parameter sets will be published as
part of the Phase~4 technical report once a candidate is chosen.

% =========================================================================
\section{Comparison to Related Systems}
\label{sec:comparison}

Several deployed and proposed systems exploit recursive SNARKs for
blockchain bootstrap or for compact verification. We summarize the
relevant ones in Table~\ref{tab:related}.

\begin{table}[h]
\centering
\small
\begin{tabularx}{\linewidth}{l X r r l}
\toprule
\textbf{System} & \textbf{Proof system} & \textbf{Proof size} & \textbf{Verify time} & \textbf{PQ?} \\
\midrule
Mina             & Pickles (Halo 2-style on Pasta)  & $\sim$200 B    & $\sim$200 ms  & No \\
Zcash (Halo 2)   & Halo 2 (Pasta)                   & $\sim$3 KB     & $\sim$10 ms   & No \\
Aleo             & Marlin (BLS12-377/BW6-761)       & $\sim$1 KB     & $\sim$20 ms   & No \\
Filecoin         & Groth16 (BLS12-381)              & 192 B          & $\sim$2 ms    & No \\
\textbf{Quillon Graph (this work)} & \textbf{Nova/BN254 $\to$ Lattice IVC} & \textbf{5--150 KB} & \textbf{5--10 ms} & \textbf{Yes (Phase 4)} \\
\bottomrule
\end{tabularx}
\caption{Recursive SNARK deployments at the time of writing. The
distinguishing property of the Quillon Graph plan is end-to-end
post-quantum security \emph{after Phase~4}, with the entire signature,
key-encapsulation, hash, and (eventually) recursion path PQ-secure.
Verify-time numbers are published benchmarks for the respective systems;
proof sizes are nominal full-finalized proof sizes.}
\label{tab:related}
\end{table}

Mina's Pickles construction is the closest precedent: a recursive SNARK
that compresses arbitrary computation, deployed in a live blockchain,
with a constant-size proof of the entire history. Mina's verify time
is slower than ours (the construction is tuned for absolute proof size
rather than verify time) and it is not post-quantum. Halo~2 and Marlin
similarly trade verify time against proof size in pairing land.
Filecoin's Groth16 is single-purpose (proof of replication / proof of
spacetime); it is not a general consensus bootstrap.

The contribution of the Quillon Graph plan, relative to these
predecessors, is the explicit post-quantum migration plan combined
with sub-10-ms verification and the operational UX integration
described in Section~\ref{sec:bootstrap}.

% =========================================================================
\section{Phased Rollout}
\label{sec:rollout}

A summary of phases is given in Table~\ref{tab:phases}. We restate
the cadence in prose here to emphasize the dependencies.

\begin{table}[h]
\centering
\small
\begin{tabularx}{\linewidth}{l l X}
\toprule
\textbf{Phase} & \textbf{Target} & \textbf{Deliverable} \\
\midrule
Phase 0  & 2026-Q2 (in flight) & Gadget library complete; SMT shadow-mode; archive-status UX live \\
Phase 1  & 2026-Q3 & $\delta$-circuit composed, single-block proof on testnet \\
Phase 2  & 2026-Q4 & Nova IVC wrapper; folded proofs on testnet \\
Phase 3  & 2027-Q3 & Advisory mainnet integration; \texttt{--bootstrap-mode=proof} opt-in \\
Phase 4  & 2028-Q1 & Lattice migration; trusted setup eliminated \\
Phase 5  & 2028-Q2/Q3 & Mandatory verification at announced activation height \\
\bottomrule
\end{tabularx}
\caption{Phased rollout. Phases 0--1 are funded engineering work
against scoped specifications; Phase 4 requires R\&D that is
scoped but not committed-pathway. Phase 5 is gated on at least
six months of advisory-mode soak with zero soundness discrepancies.}
\label{tab:phases}
\end{table}

\paragraph{Phase 0 dependencies.} Phase 0 is the only phase that is
partially complete today. The gadget library (BLAKE3, Poseidon, NTT,
Dilithium5) ships in \texttt{crates/q-ivc/}; the SMT ships in
\texttt{crates/q-storage/src/balance\_smt.rs}; the archive-status
endpoint ships in \texttt{crates/q-api-server/src/handlers.rs}. What
remains in Phase 0 is the Merkle-path gadget over BLAKE3
($\sim$600 LOC, three engineering weeks); a small range-check
wrapper around \texttt{enforce\_norm\_bound} ($\sim$3 days); and
three test-fixture corrections in the Dilithium gadget tests
($\sim$1.5 days).

\paragraph{Phase 1 dependencies.} The $\delta$-circuit composition
requires Phase~0 complete plus an additional $\sim$1{,}500 LOC of
glue connecting the existing gadgets through the per-transaction loop
and the final state-root equality check.

\paragraph{Phase 2 dependencies.} The Nova wrapper requires Phase~1
complete, plus selection between the \texttt{nova-snark} (Microsoft)
and the \texttt{arkworks-rs/nova} (community) crate. We expect to
spend approximately three engineering days prototyping both on a toy
step circuit before committing.

\paragraph{Phase 3 dependencies.} Integration into the production API
server, gossipsub topic registration, opt-in CLI flag, and a six-month
advisory soak before any further consensus-affecting decision can be
made.

\paragraph{Phase 4 dependencies.} A live, instrumented, advisory-mode
Nova deployment from Phase 3 provides the benchmark data needed to
choose among the three lattice candidates. We do not bind to a
candidate before that data is collected.

\paragraph{Phase 5 dependencies.} Six months of clean advisory
soak. A publicly announced activation height at least 20{,}000 blocks
($\sim$5.5 hours, but in practice we will announce on the order of
weeks) in the future. Old binaries continue to validate via the
block-by-block path until they choose to upgrade.

% =========================================================================
\section{Implementation Status}
\label{sec:status}

We close with an honest inventory of what is shipped, what is drafted,
and what is not yet started.

\paragraph{Shipped, in production (shadow mode where applicable):}
\begin{itemize}[leftmargin=*]
  \item \texttt{crates/q-storage/src/balance\_smt.rs} (743 LOC).
        Depth-256 sparse Merkle tree with BLAKE3 leaf and node hashes,
        empty-subtree precompute, RocksDB-backed persistence, atomic
        batch updates. Twelve unit tests covering empty tree,
        single-leaf insertion, random-tree updates, batch-vs-sequential
        equivalence, and adversarial sibling tampering.
  \item \texttt{crates/q-ivc/src/gadgets/blake3.rs} (551 LOC).
        Single-block BLAKE3 compression gadget with the
        \texttt{verify\_hash} entry point. Multi-block extension is the
        ``Path A'' handoff to an external implementer and is the
        principal blocker for the Merkle-path gadget at depth 256.
  \item \texttt{crates/q-ivc/src/gadgets/poseidon.rs} (329 LOC).
        Poseidon hash gadget; reserved for a potential v3 state
        commitment.
  \item \texttt{crates/q-ivc/src/gadgets/ntt.rs} (769 LOC).
        Cooley--Tukey decimation-in-time NTT butterfly. The
        $\mathrm{roots}[m+i] = \omega^{i \cdot n / (2m)}$ negacyclic
        convention compatible with FIPS~204 is implemented and
        tested; the BLAKE3 \texttt{FpVar}$\leftrightarrow$\texttt{UInt32}
        bridge is in place.
  \item \texttt{crates/q-ivc/src/gadgets/dilithium.rs} (902 LOC).
        Dilithium5 verification primitives:
        \texttt{compute\_az\_minus\_ct} in both standard cyclic and
        negacyclic forms, \texttt{enforce\_norm\_bound},
        \texttt{enforce\_signed\_norm\_bound} (positive-only with
        documented CAVEAT pending an upstream arkworks
        \texttt{FpVar::is\_cmp} fix), \texttt{high\_bits} per
        FIPS~204 \S5.4, and \texttt{use\_hint} per FIPS~204 \S6.5.2.
  \item \texttt{crates/q-zk-stark/} (4{,}469 LOC across nine source files).
        Full zk-STARK toolkit: \texttt{air.rs} (429 LOC) for AIR construction,
        \texttt{stark\_prover.rs} (442 LOC), \texttt{stark\_verifier.rs}
        (548 LOC), \texttt{batch\_prover.rs} (662 LOC), \texttt{polynomials.rs}
        (549 LOC) for FRI-style polynomial commitments,
        \texttt{blockchain\_state\_circuit.rs} (410 LOC),
        \texttt{wallet\_privacy\_stark.rs} (530 LOC, deployed),
        \texttt{performance.rs} (591 LOC) for benchmark harness, and a
        \texttt{gpu/} subtree for hardware acceleration.
        Approximately 70--80\% of the STARK surface is shipped; what is
        missing is the specific AIR for the Nova-SRS generator
        (Section~\ref{sec:stark-setup}) and the AIR for batched Dilithium
        verification (Section~\ref{sec:stark-elsewhere}).
  \item \texttt{crates/q-storage/src/encryption\_zkstark.rs} (562 LOC).
        \texttt{EncryptionKeyProof} for the transparent-setup attestation
        of the chain's encryption SRS. In production. The artifact
        \texttt{data-mainnet-genesis/encryption.zkstark} (176{,}832 bytes)
        is the proof itself, committed to genesis state.
\end{itemize}

\paragraph{Drafted, not wired in:}
\begin{itemize}[leftmargin=*]
  \item \texttt{crates/q-ivc/src/circuits/epoch\_transition.rs} (208 LOC).
        An \texttt{EpochTransitionCircuit} skeleton with
        \texttt{ValidatorSignatureInput} and
        \texttt{EpochTransitionInputs} types; the composition logic
        through the per-transaction loop is not yet written.
\end{itemize}

\paragraph{Not yet started:}
\begin{itemize}[leftmargin=*]
  \item \texttt{crates/q-ivc/src/gadgets/merkle.rs}. The Merkle-path
        gadget at depth 256 over BLAKE3. Blocked on the multi-block
        BLAKE3 extension (above).
  \item \texttt{crates/q-ivc/src/circuits/delta\_block.rs}. The full
        $\delta$-circuit composition. Scoped at $\sim$1{,}500 LOC.
  \item \texttt{crates/q-ivc/src/recursion/}. The Nova wrapper.
        Scoped at $\sim$800 LOC.
  \item \texttt{crates/q-api-server/src/handlers.rs} additions for
        \texttt{/api/v1/proof/tip} and \texttt{/api/v1/proof/verify}.
        Scoped at one engineering week.
  \item The Phase~4 lattice IVC implementation.
  \item The specific AIR for the Nova-SRS generator inside
        \texttt{crates/q-zk-stark/}, plus the corresponding generation
        ceremony tooling. Scoped at $\sim$300--500 LOC of AIR code
        atop the shipped \texttt{crates/q-zk-stark/} infrastructure.
\end{itemize}

\paragraph{Test debt explicitly tracked:} three Dilithium gadget tests
still use the pre-fix \texttt{roots[1] = -1} convention rather than the
post-fix \texttt{roots[1] = 1} convention for $n = 2$;
\texttt{test\_use\_hint\_with\_bias} has an off-by-one in its
expected value; \texttt{enforce\_signed\_norm\_bound} is positive-only
with a documented CAVEAT. None affect the soundness of the deployed
gadgets; all are book-keeping fixes scheduled for the next
\texttt{q-ivc} task pass.

% =========================================================================
\section{Conclusion}
\label{sec:conclusion}

A user joining the Quillon Graph network today waits hours, or accepts a
signed checkpoint and waits seconds with a trust caveat. The construction
described in this paper closes that gap. A fresh node fetches a
constant-size recursive proof, verifies it locally in about 5--10 ms,
and is immediately able to mine, transact, and serve current-state
queries --- with cryptographic certainty in the state root, not the
trust of a checkpoint operator. Historical blocks backfill in the
background, exposed through a deliberate UX surface
(\texttt{X-QNK-Archive-*} headers, HTTP 202 for not-yet-indexed
heights) so users see honest information about which queries are
answered locally and which are still being indexed.

The plan is phased deliberately. Phase~0 ships today: gadgets, SMT,
and progressive-archival UX. Phases~1--3 will roll the $\delta$-circuit,
the Nova fold, and the advisory mainnet integration through 2026 and
into 2027. Phase~4 swaps Nova for a lattice IVC scheme in early 2028,
eliminating the trusted setup and the last classical-only cryptographic
dependency in the consensus path. Mandatory verification at an
announced activation height is targeted for mid-to-late 2028 --- only
after at least six months of advisory soak with zero soundness
discrepancies. We are not in a rush. A multi-billion-dollar mainnet
does not need fast cryptography; it needs cryptography that has been
debugged in advisory mode for long enough that no one is surprised
when it becomes mandatory.

What we can promise today is that the foundations are in place: the
state commitment is shipped, the gadget library is shipped, the
zk-STARK toolkit (\texttt{crates/q-zk-stark/}, 4{,}469 LOC) and a
production transparent-setup proof at chain genesis are shipped, the
deployment cadence is published, and the architectural surface is
small enough that a Phase~4 swap of the proof system underneath does
not require touching any of the higher layers. The promise of an
instantly trustless bootstrap is no longer hypothetical --- it is a
finite, scoped engineering effort against a codebase that is already
mostly there.

We do not invent zk-STARK infrastructure for this work. We extend the
project's existing production STARK stack --- the encryption-SRS
proof at genesis, the wallet-privacy proofs, the blockchain-state
circuit, the batched prover, the AIR framework, the FRI-style
polynomial commitments --- to also attest the recursive SNARK's
setup. The same engineering composes to batched signature verification,
mobile light-client paths, cross-chain bridge attestations, and
activation-height proofs. Recursive SNARK and recursive STARK are
not competitors here: they are layered, complementary, and live in
the same crate hierarchy.

% =========================================================================
\begin{thebibliography}{99}
\bibitem{Valiant2008} Paul Valiant. Incrementally Verifiable Computation or Proofs of Knowledge Imply Time/Space Efficiency. \emph{TCC}, 2008.
\bibitem{KothapalliSettyTzialla2021} Abhiram Kothapalli, Srinath Setty, and Ioanna Tzialla. Nova: Recursive Zero-Knowledge Arguments from Folding Schemes. Cryptology ePrint Archive, Paper 2021/370, 2021.
\bibitem{KothapalliSetty2023HyperNova} Abhiram Kothapalli and Srinath Setty. HyperNova: Recursive arguments for customizable constraint systems. Cryptology ePrint Archive, Paper 2023/573, 2023.
\bibitem{BeullensSeiler2023LaBRADOR} Ward Beullens and Gregor Seiler. LaBRADOR: Compact Proofs for R1CS from Module-SIS. Cryptology ePrint Archive, Paper 2023/1119, 2023.
\bibitem{BonehChenTairi2024LatticeFold} Dan Boneh, Binyi Chen, and Akshay Tairi. LatticeFold: A Lattice-Based Folding Scheme and its Applications to Succinct Proof Systems. Cryptology ePrint Archive, Paper 2024/257, 2024.
\bibitem{NguyenSeiler2024Greyhound} Ngoc Khanh Nguyen and Gregor Seiler. Greyhound: Fast Polynomial Commitments from Lattices. Cryptology ePrint Archive, Paper 2024/1293, 2024.
\bibitem{BLAKE3spec} Jack O'Connor, Jean-Philippe Aumasson, Samuel Neves, and Zooko Wilcox-O'Hearn. BLAKE3: One Function, Fast Everywhere. Specification, BLAKE3 Team, 2020.
\bibitem{FIPS204} National Institute of Standards and Technology. Module-Lattice-Based Digital Signature Standard. FIPS PUB 204, 2024.
\bibitem{FIPS203} National Institute of Standards and Technology. Module-Lattice-Based Key-Encapsulation Mechanism Standard. FIPS PUB 203, 2024.
\bibitem{MinaPickles} Joseph Bonneau, Izaak Meckler, Vanishree Rao, and Evan Shapiro. Coda: Decentralized Cryptocurrency at Scale (Mina Protocol Technical Whitepaper). O(1) Labs, 2020.
\bibitem{Halo2} Sean Bowe, Jack Grigg, and Daira Hopwood. Halo Infinite: Proof-Carrying Data from Additive Polynomial Commitments. Cryptology ePrint Archive, Paper 2021/1536, 2021.
\bibitem{Marlin} Alessandro Chiesa, Yuncong Hu, Mary Maller, Pratyush Mishra, Noah Vesely, and Nicholas Ward. Marlin: Preprocessing zkSNARKs with Universal and Updatable SRS. Cryptology ePrint Archive, Paper 2019/1047, 2019.
\bibitem{Groth16} Jens Groth. On the Size of Pairing-Based Non-interactive Arguments. \emph{EUROCRYPT}, 2016.
\bibitem{Poseidon} Lorenzo Grassi, Dmitry Khovratovich, Christian Rechberger, Arnab Roy, and Markus Schofnegger. Poseidon: A New Hash Function for Zero-Knowledge Proof Systems. Cryptology ePrint Archive, Paper 2019/458, 2019.
\bibitem{BSBHR2018STARK} Eli Ben-Sasson, Iddo Bentov, Yinon Horesh, and Michael Riabzev. Scalable, transparent, and post-quantum secure computational integrity. Cryptology ePrint Archive, Paper 2018/046, 2018.
\bibitem{BSBHR2018FRI} Eli Ben-Sasson, Iddo Bentov, Yinon Horesh, and Michael Riabzev. Fast Reed--Solomon Interactive Oracle Proofs of Proximity. \emph{ICALP}, 2018. Cryptology ePrint Archive Paper 2017/134.
\bibitem{Cairo2021} Lior Goldberg, Shahar Papini, and Michael Riabzev. Cairo --- a Turing-complete STARK-friendly CPU architecture. Cryptology ePrint Archive, Paper 2021/1063, 2021.
\bibitem{Plonky2} Polygon Zero Team. Plonky2: Fast Recursive Arguments with PLONK and FRI. \emph{Polygon Zero}, 2022. Available at \texttt{github.com/0xPolygonZero/plonky2}.
\bibitem{BellareRogaway1993} Mihir Bellare and Phillip Rogaway. Random Oracles are Practical: A Paradigm for Designing Efficient Protocols. \emph{Proc. ACM CCS}, 1993.
\bibitem{Shor1994} Peter W. Shor. Algorithms for Quantum Computation: Discrete Logarithms and Factoring. \emph{Proc. 35th Annual Symposium on Foundations of Computer Science}, 1994.
\bibitem{NISTPQC2022} National Institute of Standards and Technology. Status Report on the Third Round of the NIST Post-Quantum Cryptography Standardization Process. NISTIR 8413, 2022.
\bibitem{KaramanolakisDAGKnight2022} Yonatan Sompolinsky, Shai Wyborski, and Aviv Zohar. DAG-KNIGHT: A Parameterless Generalization of the GHOSTDAG Protocol. Cryptology ePrint Archive, Paper 2022/1494, 2022.
\bibitem{NarwhalTusk2022} George Danezis, Lefteris Kokoris-Kogias, Alberto Sonnino, and Alexander Spiegelman. Narwhal and Tusk: A DAG-based Mempool and Efficient BFT Consensus. \emph{EuroSys}, 2022.
\end{thebibliography}

\end{document}
